% article1.tex — Статья 1 для InterNauka (CVI конференция)
% Компиляция: pdfLaTeX
% =====================================================================

\documentclass[14pt]{extarticle}

% --- Шрифты и язык ---
\usepackage[T2A]{fontenc}
\usepackage[utf8]{inputenc}
\usepackage[english,russian]{babel}
\usepackage{tempora}            % Times New Roman (Cyrillic)

% --- Геометрия страницы ---
\usepackage[left=2cm, right=2cm, top=2cm, bottom=2cm]{geometry}

% --- Межстрочный интервал ---
\usepackage{setspace}
\onehalfspacing

% --- Отступ первого абзаца ---
\usepackage{indentfirst}
\setlength{\parindent}{1.25cm}

% --- Таблицы ---
\usepackage{array}
\usepackage{booktabs}
\usepackage{longtable}
\usepackage{multirow}
\usepackage{caption}
\captionsetup{justification=centering, labelsep=period, font={normalsize}}

% --- Без нумерации страниц (опционально; раскомментировать при необходимости) ---
% \pagestyle{empty}

\begin{document}

% =====================================================================
% ЗАГОЛОВОК
% =====================================================================

\begin{center}
\textbf{РАЗРАБОТКА СТРУКТУРЫ ЭЛЕКТРОННОГО ОБРАЗОВАТЕЛЬНОГО РЕСУРСА ПО ФИЗИКЕ ДЛЯ ОРГАНИЗАЦИИ УСВОЕНИЯ СИСТЕМЫ ФИЗИЧЕСКИХ ЗНАНИЙ НА ОСНОВЕ КОНЦЕПЦИИ П.\,Я.\,ГАЛЬПЕРИНА}
\end{center}

\begin{flushright}
% TODO: уточнить отчество автора
\textbf{\textit{Синяков Вячеслав Андреевич}}\\
\textit{магистрант,\\
Московский Политехнический Университет,\\
РФ, г. Москва}

\vspace{0.3cm}
\textbf{\textit{Прояненкова Лидия Алексеевна}}\\
\textit{доктор педагогических наук, профессор,\\
Московский педагогический государственный университет,\\
РФ, г. Москва}
\end{flushright}

\vspace{0.5cm}

% =====================================================================
% АННОТАЦИЯ
% =====================================================================

\noindent\textbf{АННОТАЦИЯ}

\noindent
Статья посвящена проектированию и реализации электронного образовательного ресурса (ЭОР)
по физике для учащихся 7--9 классов. Ресурс построен на основе концепции поэтапного формирования
умственных действий П.\,Я.\,Гальперина и реализует три варианта последовательности учебных
действий, выделенных Л.\,А.\,Прояненковой для усвоения элементов физических знаний.
Описаны принципы построения модулей ЭОР: модуля ознакомления с системой знаний,
модуля адаптивного решения задач, модуля составления алгоритма решения, модуля пооперационного
контроля и модуля итогового контроля. Раскрыта логика адаптивного ветвления, при которой
система автоматически переводит ученика на более детализированный вариант работы в случае
затруднений. Представлены ключевые проектные решения, связанные с клиент-серверной
архитектурой, структурой базы данных и интерфейсом учителя для наполнения ресурса
содержанием.

\vspace{0.3cm}
\noindent\textbf{Ключевые слова:} электронный образовательный ресурс, физика, основная школа, теория поэтапного формирования умственных действий, П.\,Я.~Гальперин, система знаний, адаптивное обучение.

\vspace{0.5cm}

% =====================================================================
% ВВЕДЕНИЕ
% =====================================================================

\section*{Введение}

Федеральный государственный образовательный стандарт основного общего образования [9]
предъявляет к выпускникам основной школы требования, связанные с овладением
системой физических знаний: понятиями о явлениях, физическими величинами, законами.
Учащиеся должны уметь применять эти знания для анализа конкретных ситуаций.
Практика обучения физике показывает, что формирование такого умения требует
многократного выполнения учащимися действий по применению знаний в разнообразных
ситуациях~--- не менее шести--восьми раз для каждого элемента знания [6].
В условиях ограниченного учебного времени организовать такую работу на уроке затруднительно.

Электронные образовательные ресурсы (ЭОР) позволяют организовать
индивидуальную работу учащихся, автоматизировать контроль и обеспечить
оперативную обратную связь. Анализ доступных ЭОР по физике, однако,
показывает их методическую ограниченность: ресурсы
сводятся к демонстрационным материалам и подборкам задач с выбором ответа [10].
Методическая модель, заложенная в них, воспроизводит схему <<предъявление
информации~--- проверка запоминания>>. ЭОР, который организует поэтапное
формирование действий по применению физических знаний, в школьной практике
не представлен.

Теоретической основой для проектирования подобного ресурса служит концепция
поэтапного формирования умственных действий П.\,Я.\,Гальперина [3].
Л.\,А.\,Прояненкова обосновала возможность создания ЭОР по физике на основе
этой концепции, выделила три варианта последовательности учебных действий
в зависимости от степени трудности осваиваемого материала для ученика
и предложила блок-схему обучающей программы [6]. Программная реализация
описанного ресурса до настоящего времени не была осуществлена.

Цель статьи~--- описать разработанную структуру ЭОР по физике для основной
школы, реализующую три варианта методики обучения на основе концепции
П.\,Я.\,Гальперина, и обосновать принятые проектные решения.


% =====================================================================
% 1. ТЕОРЕТИЧЕСКИЕ ОСНОВАНИЯ
% =====================================================================

\section{Теоретические основания проектирования ЭОР}

\subsection{Концепция П.\,Я.\,Гальперина}

Концепция поэтапного формирования умственных действий и понятий (ТПФУД)
рассматривает усвоение знаний как процесс интериоризации~--- перехода действия
из внешней формы во внутреннюю через ряд закономерных этапов.
П.\,Я.\,Гальперин выделил шесть таких этапов: мотивационный; составление
схемы ориентировочной основы действия (ООД); выполнение действия
в материализованной форме; в громкой речи; во внешней речи <<про себя>>;
во внутренней речи [3].

Ключевое положение концепции состоит в том, что качество сформированного действия
определяется содержанием и характером ориентировочной основы, которую получает
обучаемый [5, 8]. Если ООД содержит указания на существенные признаки и отношения
в обобщённой форме, формируемое действие приобретает свойства разумности,
обобщённости и сознательности.

Применительно к основной школе этапы громкой речи и внешней речи <<про себя>>
объединяются в один этап, что даёт четыре рабочих этапа:
мотивационный, составление схемы ООД, выполнение действия (в материализованной,
речевой и умственной формах) и контроль [6].


\subsection{Система знаний и действия по её применению}

В рамках деятельностного подхода к обучению физике единицей анализа учебного
процесса является действие, адекватное усваиваемому элементу знания [1, 4].
Для каждого вида физических знаний (понятие о явлении, физическая величина,
закон) выделены логические схемы деятельности по получению и применению
этих знаний [1, 7].

Система знаний о физическом явлении включает: название явления, определение,
условия протекания, характеристики (физические величины), законы,
описывающие явление, и практические применения.
Типовая задача состоит в том, чтобы определить значение одной из величин,
описывающих явление, в конкретной ситуации. Метод решения типовой задачи
представляет собой последовательность действий: выделение движущегося тела,
определение начального положения и характеристик, выделение участков движения,
составление уравнений для каждого участка и т.\,д. Учащийся, освоивший эту
последовательность, может применить её к различным ситуациям.


\subsection{Три варианта последовательности учебных действий}

Полнота выполнения и последовательность учебных действий зависят от того,
насколько трудным для конкретного ученика оказывается осваиваемый материал.
Л.\,А.\,Прояненкова выделила три варианта работы, соответствующие разной
степени трудности [6]. Они представлены в таблице~1.

\begin{table}[htbp]
\centering
\caption{Варианты последовательности действий ученика по усвоению элемента знания}
\label{tab:variants}
\small
\begin{tabular}{|p{3.5cm}|p{5.5cm}|c|c|c|}
\hline
\textbf{Этап усвоения} & \textbf{Действие учащегося} & \textbf{В\,1} & \textbf{В\,2} & \textbf{В\,3} \\
\hline
\multirow{3}{3.5cm}{1. Мотивационный}
  & 1.1. Выбор элемента знания    & + & + & + \\
  & 1.2. Выбор действия            & + & + & + \\
  & 1.3. Знакомство с~заданием     & + & + & + \\
\hline
2. Составление схемы ООД
  & 2. Составление алгоритма       & -- & + & + \\
\hline
\multirow{3}{3.5cm}{3. Выполнение действия}
  & 3а. По алгоритму с~контролем каждой операции & -- & -- & + \\
  & 3б. С~объяснением операций      & -- & -- & + \\
  & 3в. Самостоятельно              & + & + & + \\
\hline
4. Контроль
  & 4. Задание на~оценку            & + & + & + \\
\hline
\end{tabular}
\end{table}

Вариант~1 рассчитан на ученика, для которого задание не представляет
трудности: после ознакомления он выполняет его самостоятельно и переходит
к контрольному заданию. Вариант~2 предполагает, что ученик испытывает
затруднения; ему предлагается составить алгоритм выполнения задания,
опираясь на метод решения типовой задачи, после чего он решает задачу
самостоятельно. Вариант~3 предназначен для случаев значительных затруднений:
ученик проходит через материализованную и речевую формы выполнения действия
с пооперационным контролем каждого шага.

Переход между вариантами определяется результатами работы.
Ученик, начавший по варианту~1, при ошибках переводится
на вариант~2 или~3. Управление учебной работой в зависимости
от варианта~--- задача, которую Л.\,А.\,Прояненкова предложила
передать компьютеру [6]. Первая попытка создания соответствующего
ЭОР <<Учимся применять физические знания>> была предпринята
в 2018~году на платформе WPF (C\#) [6]; настоящая работа развивает
этот замысел в форме веб-приложения с полным циклом адаптивного
управления.


% =====================================================================
% 2. СТРУКТУРА ЭОР
% =====================================================================

\section{Структура электронного образовательного ресурса}

Структура ЭОР определяется логикой трёх вариантов учебной работы.
Ресурс состоит из пяти модулей, каждый из которых соответствует
определённому этапу усвоения знаний. Общая последовательность
взаимодействия ученика с ресурсом выглядит следующим образом:
ознакомление с системой знаний, оценка трудности, решение задач
с адаптивным ветвлением, итоговый контроль.

\textbf{Модуль ознакомления с системой знаний} соответствует
мотивационному этапу.
Ученику предъявляется таблица-система знаний о физическом явлении~---
например, <<Равномерное и неравномерное движение на участке траектории>>.
В таблице сведены название явления, определение, характеристики, формулы
и единицы измерения. Далее ученик знакомится с условиями всех задач:
видит перечень ситуаций, в которых ему предстоит применить систему знаний.
Решать задачи на этом этапе не требуется~--- ученик составляет общее
представление о предстоящей работе.

После ознакомления ученик оценивает степень трудности задания для себя.
Ресурс предлагает выбор из трёх уровней:
<<нетрудно>> (вариант~1), <<трудно>> (вариант~2) и <<очень трудно>> (вариант~3).
Рядом с каждым уровнем приведено пояснение, какой объём помощи
получит ученик при данном выборе. По результатам самооценки система
определяет начальный вариант работы.

\textbf{Модуль решения задач}~--- центральный компонент ЭОР,
обеспечивающий выполнение действия в умственной форме (этап~3в).
Ученику предъявляется задача из банка, привязанного к данной системе знаний.
Решение выполняется в тетради; в ресурс вводится численный ответ,
при желании~--- фотография записей. Проверка ответа учитывает
допустимую погрешность и единицы измерения.

Адаптивный алгоритм выбора следующей задачи учитывает историю ответов ученика.
Базовый уровень сложности определяется самооценкой трудности; если ученик
решает задачи с долей правильных ответов выше 80\,\%, уровень повышается;
если доля правильных ответов падает ниже 60\,\% или ученик допускает два
подряд неправильных ответа, уровень понижается.

Адаптивное ветвление по ошибкам работает в рамках текущей задачи.
Если ученик дважды ответил неверно, ресурс предлагает перейти к составлению
алгоритма решения (модуль~3). Если ученик трижды ответил неверно, ресурс
переводит его к пооперационному контролю (модуль~4). Такая логика
соответствует переходу от варианта~1 к вариантам~2 и~3 в таблице~1.

В \textbf{модуле составления алгоритма решения} (этап~2: составление
схемы ООД) ученик работает с текстом метода решения типовой задачи.
В тексте ряд ключевых слов заменён пропусками. Рядом расположено
<<облако слов>>~--- перемешанный набор правильных ответов и дистракторов.

Ученик заполняет пропуски, перетаскивая слова из облака на место пропуска
или выбирая их нажатием. Использованные слова исчезают из облака.
Если одно и то же слово встречается в нескольких пропусках, рядом
с ним отображается счётчик оставшихся использований.
После заполнения всех пропусков ученик нажимает кнопку <<Проверить>>;
система сопоставляет его ответы с эталоном и указывает
на ошибочно заполненные позиции.

Содержание этого модуля формируется учителем.
В базе данных хранятся шаги метода решения: каждый шаг имеет заголовок
(например, <<Выделите движущееся тело>>) и описание. Учитель отмечает слова,
которые станут пропусками, и при необходимости добавляет дистракторы.

\textbf{Модуль пооперационного контроля} переводит работу ученика
в материализованную форму (этап~3а). Задача разбивается на операции
по шагам метода решения. Ученик вводит промежуточный результат
каждой операции и получает немедленную обратную связь. Пошаговый
режим позволяет локализовать конкретное действие, в котором
допущена ошибка.

\textbf{Модуль итогового контроля} (этап~4) отличается от предыдущих
отсутствием обратной связи. Ученик получает задание <<на оценку>>, вводит
ответ и загружает фотографию решения. Результат проверяет учитель.
Условия приближены к реальному контролю: ученик вынужден самостоятельно
проверить решение до отправки.


% =====================================================================
% 3. ПРОГРАММНАЯ РЕАЛИЗАЦИЯ
% =====================================================================

\section{Программная реализация и интерфейс}

ЭОР реализован как веб-приложение. Серверная часть написана на
языке Python с применением фреймворка Django и расширения
Django REST Framework, предоставляющего программный интерфейс (API).
Клиентская часть использует библиотеку React (JavaScript);
данные хранятся в СУБД PostgreSQL. Веб-технологии выбраны
для того, чтобы обеспечить работу ресурса в браузере на любом
устройстве~--- компьютере, планшете, смартфоне~--- без установки
дополнительных программ.

Структура базы данных отражает предметную область.
Модель <<Система знаний>> (KnowledgeSystem) содержит название,
описание и изображение таблицы-системы знаний; с ней связаны задачи (Task),
метод решения типовой задачи (SolutionMethod) и шаги этого метода (SolutionStep).
Каждая задача имеет текст условия, правильный ответ, единицу измерения,
допустимую погрешность и уровень сложности (от~1 до~5). Шаги метода решения
хранят заголовок, описание и порядковый номер. Для организации cloze-заданий
(заполнения пропусков) предусмотрены поля, содержащие исходный текст,
размеченный текст с отмеченными пропусками и список дистракторов.

Модель учебной сессии (LearningSession) фиксирует состояние работы
ученика: текущий этап, выбранный вариант трудности, количество решённых
и правильно решённых задач, число ошибок подряд. Попытки решения задач
(TaskAttempt) и попытки выполнения отдельных операций (StepAttempt)
сохраняются для анализа учебных результатов.

Интерфейс учителя позволяет создавать и редактировать системы знаний,
задачи и методы решения. Учитель формирует cloze-задания:
вводит текст алгоритма решения, отмечает слова для пропуска,
добавляет дистракторы. Для ускорения работы предусмотрена
функция автоматического формирования текста cloze-задания из шагов
метода решения, хранящихся в базе данных.

Интерфейс ученика организован как последовательность экранов.
Ученик видит только текущий этап; информация об алгоритме
адаптивного управления от него скрыта~--- переходы между вариантами
происходят автоматически.
Экран ознакомления с задачами разделён на две колонки: слева~---
прокручиваемый список ситуаций с условиями, справа~--- кнопка
перехода к работе и пояснения. Такой макет исключает необходимость
пролистывать длинный перечень задач до управляющих элементов.
Шрифты, цвета и размеры элементов подобраны
с учётом возрастной группы 7--9 классов.


% =====================================================================
% ЗАКЛЮЧЕНИЕ
% =====================================================================

\section*{Заключение}

Описана структура электронного образовательного ресурса по физике для
учащихся 7--9 классов, спроектированного на основе концепции поэтапного
формирования умственных действий П.\,Я.\,Гальперина. Ресурс реализует
три варианта последовательности учебных действий, предложенные
Л.\,А.\,Прояненковой [6], и обеспечивает автоматическое управление
переходами между ними в зависимости от результатов работы ученика.

Пять модулей~--- ознакомления с системой знаний, решения задач,
составления алгоритма, пооперационного контроля и итогового контроля~---
покрывают этапы усвоения от мотивации до контроля.
Интерфейс учителя позволяет наполнять ресурс содержанием
по любой теме школьного курса физики, не прибегая к программированию.

Следующий этап работы~--- апробация ресурса в учебном процессе
основной школы. Задачи апробации: проверить, приводит ли применение
ЭОР к повышению качества усвоения систем физических знаний по сравнению
с традиционным обучением, и уточнить пороги адаптивного
алгоритма (число допустимых ошибок, критерии повышения сложности).
Банк заданий планируется расширить за пределы кинематики на темы
<<Тепловые явления>>, <<Электрические явления>>, <<Световые явления>>.


% =====================================================================
% СПИСОК ЛИТЕРАТУРЫ
% =====================================================================

\section*{Список литературы}

\begin{enumerate}
\item Анофрикова С.\,В., Прояненкова Л.\,А. Руководство по разработке
  уроков с использованием учебного физического эксперимента: учеб. пособие
  для студентов вузов. Астрахань: Астрах. ун-т, 2005. 78~с.

\item Анофрикова С.\,В., Стефанова Г.\,П. Практическая методика преподавания
  физики. Ч.~1: Учебное пособие. Астрахань: Изд-во Астраханского пед. ин-та,
  1995. 232~с.

\item Гальперин П.\,Я. Методы обучения и умственное развитие ребенка.
  М.: Изд-во МГУ, 1985. 45~с.

\item Прояненкова Л.\,А. Деятельностный подход в обучении физике //
  Физика в школе. 2005. №~1. С.~34--41.

\item Прояненкова Л.\,А. Теория поэтапного формирования умственных действий
  в методической подготовке учителя физики // Преподаватель XXI ВЕК. 2009.
  №~3. С.~59--66.

% TODO: уточнить название сборника и год издания у научного руководителя
\item Прояненкова Л.\,А. ЭОР по физике на основе концепции П.\,Я.\,Гальперина //
  Физика и её преподавание в школе и вузе: сб. материалов конф.
  [уточнить выходные данные]. С.~137--141.

\item Прояненкова Л.\,А., Нгуен Ву~Ань. Методика формирования знания о логике
  экспериментального исследования и экспериментальных умений у учащихся средней
  школы Вьетнама при изучении физики // Школа будущего. 2020. №~4.
  С.~250--267.

\item Талызина Н.\,Ф. Педагогическая психология: учеб. пос. для студ. сред. пед.
  учеб. завед. М.: Академия, 1998. 288~с.

\item Федеральный государственный образовательный стандарт основного общего
  образования: утв. приказом Министерства просвещения Российской Федерации
  от 31.05.2021 №~287.

\item Якушева Н.\,М. Дидактические принципы создания средств электронного
  обучения и вопросы их реализации // Вестник университета (ГУУ). 2011.
  №~16. С.~49--55.
\end{enumerate}

\end{document}
